\RequirePackage{fix-cm}
\documentclass[twocolumn]{svjour3}

\usepackage{cite}
\usepackage{graphicx}
\usepackage{amsmath}
\usepackage{amsfonts}
\usepackage{amssymb}
\usepackage{xcolor}
\usepackage{mathtools}
\usepackage{bm}
\usepackage{upgreek}
\usepackage{adjustbox}
\usepackage{placeins}    
\usepackage{balance}
\usepackage[misc]{ifsym}
\usepackage{url}
\usepackage{multirow}
\usepackage{tikz}
\usepackage{soul}

\newcommand\hhl[1]{%
	\bgroup
	\hskip0pt\color{blue!80!blue}%
	#1%
	\egroup
}

%%%%%%%%%%%%%%%%%%%%%%%%%%%

\journalname{Nonlinear Dynamics}
	
\begin{document}
	
	\title{\hl{Nie uzupelniac}}
	
	\titlerunning{\hl{Nie uzupelniac}}
	
	\author{Adam~Soko{\l}owski, Robert~Nebeluk}
	
	\authorrunning{A.~Soko{\l}owski, R.~Nebeluk} 
	
	\institute{A.~Soko{\l}owski, R.~Nebeluk, \at
		Warsaw University of Technology\\
		Institute of Control and Computation Engineering\\  
		ul. Nowowiejska 15/19, 00-665 Warsaw, Poland\\
		\email{ \hl{a.sokolowski2003@gmail.com}, Robert.Nebeluk@pw.edu.pl}
	}
	
	\date{Received: date / Accepted: date}
	% The correct dates will be entered by the editor

	\maketitle
	
	\begin{abstract}
		\hl{Nie uzupelniac}
		\keywords{\hl{Nie uzupelniac}}
	
	\end{abstract}
	
	\section{Introduction}
	\label{sec:p}
    \hl{Nie uzupelniac}\\
    \hl{Na wszelki wypadek: prosze pamietac aby nie generowac tekstu przy pomocy narzedzi AI - pliki zrodlowe po wyslaniu beda sprawdzane przez odpowiednie narzedzia. Dopoki sam Pan tlumaczy i sam Pan wymysla kolejne linijki tekstu to na pewno bedzie wszystko w porzadku.}
    \section{Modeling of inverted pendulum on cart system}
	\label{sec:pp}
    \hl{TODO: Wykorzystac material z pracy inżynierskiej: rozdzial 2 (bez rozdz. 2.5.1, 2.5.2, 2.6) oraz tabelka 3.1}
The inverted pendulum on a cart is an underactuated mechanical system that serves as a
standard benchmark for evaluating control strategies. The system
consists of a cart of mass $M$ constrained to horizontal translation, with a rigid pendulum of
mass $m$ and length $l$ pinned at the cart centre. The sole actuator applies a horizontal force
$u(t)$ to the cart; an external horizontal disturbance $F_\mathrm{w}(t)$ acting on the pendulum
tip is also considered. The physical parameters employed throughout this work are listed in
Table~\ref{tab:params}.

\setlength\tabcolsep{4pt}
\begin{table}[htbp]
\caption{Physical parameters of the inverted-pendulum system.}
\label{tab:params}
\centering
\begin{small}
\begin{tabular}{lcc}
\hline
\rule{0pt}{3mm}Parameter & Symbol & Value \\
\hline
\rule{0pt}{3.5mm}Cart mass         & $M$ & $2.40$~kg \\
Pendulum mass     & $m$ & $0.23$~kg \\
Pendulum length   & $l$ & $0.36$~m  \\
Gravitational acceleration & $g$ & $9.81$~m/s$^{2}$ \\
\hline
\end{tabular}
\end{small}
\end{table}

The following simplifying assumptions are adopted: the rod is massless and
rigid; friction at the wheels and at the pivot is negligible; motion is planar; and the
pendulum mass is treated as a point mass at the rod tip. The coordinate $\theta(t)$ denotes the
pendulum deflection from the upright equilibrium (positive clockwise, $\theta=0$ at the top),
and $x(t)$ is the cart displacement.

\subsection{Equations of motion}

The position of the pendulum tip in the fixed frame is
$x_G = x + l\sin\theta$, $y_G = l\cos\theta$.
Differentiating twice and applying Newton's second law to the whole system in the horizontal
direction gives:
\begin{equation}
(M{+}m)\ddot{x} + ml\cos\theta\,\ddot{\theta} - ml\sin\theta\,\dot{\theta}^{2} = u + F_\mathrm{w}.
\label{eq:trans}
\end{equation}
Balancing moments about the pivot yields:
\begin{equation}
ml\cos\theta\,\ddot{x} + ml^{2}\ddot{\theta} = mgl\sin\theta + lF_\mathrm{w}\cos\theta.
\label{eq:rot}
\end{equation}
Introducing $M_T = M+m$ and $D(\theta)=l\bigl(M_T - m\cos^{2}\!\theta\bigr)$, the explicit
accelerations resolved from \eqref{eq:trans}--\eqref{eq:rot} are:
\begin{align}
\ddot{x} &= \frac{u + ml\sin\theta\,\dot{\theta}^{2} - mg\sin\theta\cos\theta}{M_T - m\cos^{2}\!\theta},
\label{eq:xdd}\\[2pt]
\ddot{\theta} &= \frac{M_T g\sin\theta - \cos\theta\!\left(u + ml\sin\theta\,\dot{\theta}^{2}\right)}{D(\theta)}.
\label{eq:tdd}
\end{align}

\subsection{Nonlinear state-space model}

Defining the state vector $x(t) = [\theta,\dot{\theta},x,\dot{x}]^T$,
the dynamics \eqref{eq:xdd}--\eqref{eq:tdd} take the form $\dot{x} = f(x,u,F_\mathrm{w})$:
\begin{equation}
\begin{bmatrix}\dot{x}_1\\\dot{x}_2\\\dot{x}_3\\\dot{x}_4\end{bmatrix}
=
\begin{bmatrix}
x_2\\[2pt]
\dfrac{M_T g\sin x_1 - \cos x_1\!\left(u + ml\sin x_1\,x_2^{2}\right)}{D(x_1)}\\[6pt]
x_4\\[2pt]
\dfrac{u + ml\sin x_1\,x_2^{2} - mg\sin x_1\cos x_1}{M_T - m\cos^{2}\!x_1}
\end{bmatrix}.
\label{eq:nonlin_ss}
\end{equation}

\subsection{Linearized model}

For small deflections around the upright equilibrium ($\theta\approx 0$), the approximations
$\sin\theta\approx\theta$, $\cos\theta\approx 1$ yield the linear time-invariant
model $\dot{x} = Ax + Bu + B_\mathrm{d}F_\mathrm{w}$:
\begin{equation}
A = \begin{bmatrix}0&1&0&0\\ \dfrac{M_T g}{Ml}&0&0&0\\[4pt]0&0&0&1\\[2pt]-\dfrac{mg}{M}&0&0&0\end{bmatrix},\;
B = \begin{bmatrix}0\\[2pt]-\dfrac{1}{Ml}\\[4pt]0\\[2pt]\dfrac{1}{M}\end{bmatrix}.
\label{eq:linear}
\end{equation}
The eigenvalues of $A$ are $\lambda_{1,2}=0$ and $\lambda_{3,4}\approx\pm 5.49$; the positive
real eigenvalue confirms open-loop instability. The pair $(A,B)$ has full controllability rank
equal to four, guaranteeing the existence of a stabilising feedback law.
    
    \section{Control system design}
	\label{sec:r}
	\hl{TODO: Wykorzystac material z pracy inżynierskiej: rozdzial 3 (tylko opis matematyczny/implementacyjny - bez rysunkow/wynikow/strojenia)}

 
All controllers operate in discrete time with sampling period $T_p$ and share the state vector
$$x(k) = [\theta(k),\dot{\theta}(k),x(k),\dot{x}(k)]^T$$
The control output $u(k)$ is the horizontal force applied to the cart at time instant
$t = k T_p$.

\subsection{PID control scheme}

The parallel PID scheme employs two independent PID loops whose outputs are summed:
\begin{equation}
u(k) = u_{\theta}(k) + u_{x}(k),
\end{equation}
where $e_{\theta}(k)=\theta_\mathrm{ref}-x_1(k)$ and $e_x(k)=x_\mathrm{ref}-x_3(k)$.
Each loop follows a discrete-time PID law with \emph{derivative on measurement} to suppress
quantisation-noise amplification:
\begin{align}
u_{\theta}(k) &= K_{p,\theta}e_{\theta}(k) + K_{i,\theta}\!\sum_{j=0}^{k}\!e_{\theta}(j)T_p
               - K_{d,\theta}x_2(k),\label{eq:pid_theta}\\
u_{x}(k) &= K_{p,x}e_x(k) + K_{i,x}\!\sum_{j=0}^{k}\!e_x(j)T_p
           - K_{d,x}x_4(k).\label{eq:pid_x}
\end{align}
Using the velocity states $x_2$ and $x_4$ directly as derivative feedback eliminates
derivative kick during reference steps. An anti-windup mechanism limits the integral
accumulation to $\pm 5.0$.

\subsection{LQR control scheme}

The Discrete Linear Quadratic Regulator (DLQR)~\cite{Jezierski2017} finds the
state-feedback gain $K$ minimising the infinite-horizon cost:
\begin{equation}
J = \sum_{k=0}^{\infty}\!\left(x(k)^T Q\,x(k) + u(k)^T R\,u(k)\right),
\label{eq:lqr_cost}
\end{equation}
subject to the linearised discrete model $x(k{+}1)=A_dx(k)+B_du(k)$,
where $A_d,B_d$ are obtained from $(A,B)$ of \eqref{eq:linear} via zero-order-hold (ZOH)
discretisation at period $T_p$. The gain is computed from the unique positive-definite
solution $P$ of the Discrete Algebraic Riccati Equation (DARE):
\begin{equation}
P = A_d^T P A_d - A_d^T P B_d(R+B_d^T P B_d)^{-1}B_d^T P A_d + Q,
\end{equation}
yielding:
\begin{equation}
K = (R + B_d^T P B_d)^{-1}B_d^T P A_d,
\end{equation}
and the control law $u(k) = -K\bigl(x(k) - x_\mathrm{ref}\bigr)$,
where $Q\succeq 0$, $R>0$ are design weighting matrices.

\subsection{LQR-I control scheme}

To eliminate steady-state cart-position error, the DLQR is augmented with an integrator
acting on the position tracking error. The accumulated error state is:
\begin{equation}
z(k+1) = z(k) + \bigl(x_\mathrm{ref} - x_3(k)\bigr)T_p.
\end{equation}
The augmented state $x_\mathrm{aug}(k)=[x_1,x_2,x_3,x_4,z]^T$ obeys:
\begin{equation}
x_\mathrm{aug}(k{+}1) = \tilde{A}_d\,x_\mathrm{aug}(k) + \tilde{B}_d\,u(k),
\end{equation}
where $\tilde{A}_d\in\mathbb{R}^{5\times 5}$ and $\tilde{B}_d\in\mathbb{R}^{5\times 1}$
append the integrator row to $A_d$ and $B_d$, respectively. Applying DLQR to the augmented
system with extended weighting matrices $\tilde{Q}$ and $R$ yields:
\begin{equation}
u(k) = -\tilde{K}\,x_\mathrm{aug}(k),
\end{equation}
providing zero steady-state position error for constant references.

\subsection{PID-LQR control scheme}

The PID-LQR scheme combines DLQR stabilisation with an outer PID loop on the position
error:
\begin{equation}
u(k) = u_\mathrm{LQR}(k) + u_\mathrm{PID,pos}(k),
\end{equation}
where the LQR component acts on the full state:
\begin{equation}
u_\mathrm{LQR}(k) = -K\bigl(x(k) - x_\mathrm{ref}\bigr),
\end{equation}
and $u_\mathrm{PID,pos}(k)$ follows the discrete PID law of \eqref{eq:pid_x} restricted to the
cart-position loop. This additive structure improves step-reference tracking and rejects
steady-state errors due to model uncertainty.

\subsection{Fuzzy LQR control scheme}

The Fuzzy-LQR scheme operates a Takagi-Sugeno (T-S) fuzzy block in parallel with the
DLQR:
\begin{equation}
u(k) = u_\mathrm{LQR}(k) + G\cdot u_\mathrm{Fuzzy}(k).
\end{equation}
Triangular membership functions over two linguistic sets (\emph{small error},
\emph{large error}) for each of the four state variables yield a rule base of
$2^4=16$ rules of the form:
\begin{quote}
\textbf{IF} $e_\theta$ is $A_i$ \textbf{AND} $\dot\theta$ is $B_i$
$\ldots$ \textbf{THEN} $u_i = -K^{(i)}x$.
\end{quote}
The local gain vectors adapt to the severity of the deviation:
\begin{equation}
K^{(i)} = K_\mathrm{base} + \alpha\,L_i,
\end{equation}
where $L_i\in\{0,\ldots,4\}$ is the number of state variables currently in the
\emph{large-error} region, and $\alpha$ is a fixed increment. The defuzzified
output is the firing-strength-weighted average:
\begin{equation}
u_\mathrm{Fuzzy}(k) = \frac{\displaystyle\sum_{i=1}^{16}w_i(x)\,u_i}
                           {\displaystyle\sum_{i=1}^{16}w_i(x)},
\end{equation}
where $G$ is a global scaling gain. This adaptive-gain structure intensifies the
corrective action at large deflections while remaining benign near the equilibrium.

\subsection{MPC-NO control scheme}

The Nonlinear Model Predictive Controller (NMPC) solves, at each sampling instant, the
finite-horizon optimisation~\cite{Camacho2007,Mills2009}:
\begin{equation}
\min_{\Delta\u(k)}\;J(k) = \sum_{p=1}^{N}e_p^T Q\,e_p
                              + \sum_{p=0}^{N_u-1}\Delta u_p^T R\,\Delta u_p,
\label{eq:mpc_cost}
\end{equation}
where $e_p(k)=x_\mathrm{ref}(k{+}p|k)-\hat{x}(k{+}p|k)$,
$\Delta u_p = u(k{+}p|k)-u(k{+}p{-}1|k)$, $N$ is the prediction horizon and
$N_u\leq N$ the control horizon, subject to:
\begin{align}
\hat{x}_{k+1} &= f\!\left(\hat{x}_k,u_k\right),\quad k=0,\ldots,N{-}1,\\
u_{\min} &\leq u_k \leq u_{\max}.
\end{align}
The prediction model $f(\cdot)$ integrates the nonlinear dynamics \eqref{eq:nonlin_ss} over
one sampling period using fourth-order Runge-Kutta (RK4) with a ZOH assumption on the input.
The optimisation is solved by the SLSQP algorithm from \texttt{scipy.optimize}.

\subsection{MPC-NO-alt control scheme}

The MPC-NO-alt variant extends the cost \eqref{eq:mpc_cost} with a direct penalty on control
amplitude to reduce actuator wear and energy consumption:
\begin{equation}
J(k) = \sum_{p=1}^{N}e_p^T Q\,e_p
      + \sum_{p=0}^{N_u-1}\!\left(\Delta u_p^T R\,\Delta u_p + u_p^T R_\mathrm{abs}\,u_p\right),
\end{equation}
where $R_\mathrm{abs}>0$ penalises the absolute force magnitude. The prediction model,
constraint structure, and solver are identical to the MPC-NO formulation.

\subsection{LMPC control scheme}

The Linear MPC (LMPC) replaces the nonlinear RK4 prediction model with the ZOH-discretised
linear model obtained from \eqref{eq:linear}:
\begin{equation}
\hat{x}_{k+1} = A_d\hat{x}_k + B_du_k,\quad k=0,\ldots,N{-}1,
\end{equation}
reducing prediction to fast matrix recursions. The cost function is identical to
\eqref{eq:mpc_cost} and the same SLSQP solver is used. Although the linear model introduces
approximation errors at large deflections, it substantially lowers the per-step computational
burden compared with NMPC.

\section{Control quality indicators}
\hl{TODO: Wykorzystac material z pracy inżynierskiej: rozdzial 5.2}
To provide a quantitative and objective comparison of the evaluated control strategies, the
following performance indices are computed over $N$ discrete samples of the controlled
output $y(k)$ (pendulum angle $\theta$ or cart position $x$):
\begin{itemize}
\item \textbf{Mean Squared Error} (MSE) --- penalises large deviations strongly:
\begin{equation}
\mathrm{MSE} = \frac{1}{N}\sum_{k=1}^{N}\bigl(y(k)-y_\mathrm{ref}(k)\bigr)^2.
\end{equation}
\item \textbf{Mean Absolute Error} (MAE) --- average absolute deviation:
\begin{equation}
\mathrm{MAE} = \frac{1}{N}\sum_{k=1}^{N}\bigl|y(k)-y_\mathrm{ref}(k)\bigr|.
\end{equation}
\item \textbf{Integral of Squared Error} (ISE) --- cumulative quadratic error energy:
\begin{equation}
\mathrm{ISE} = \sum_{k=1}^{N}\bigl(y(k)-y_\mathrm{ref}(k)\bigr)^2.
\end{equation}
\item \textbf{Integral of Absolute Error} (IAE) --- cumulative absolute error:
\begin{equation}
\mathrm{IAE} = \sum_{k=1}^{N}\bigl|y(k)-y_\mathrm{ref}(k)\bigr|.
\end{equation}
\item \textbf{Quadratic control cost} $E_{u,L_2}$ and \textbf{absolute control cost}
$E_{u,L_1}$, measuring actuator effort:
\begin{align}
E_{u,L_2} &= \sum_{k=1}^{N}u(k)^{2},\\
E_{u,L_1} &= \sum_{k=1}^{N}|u(k)|.
\end{align}
\item \textbf{Settling time} $T_s$ --- the first instant after which
$|\theta(k)|<\varepsilon_\theta = 0.001$~rad and $|x(k)|<\varepsilon_x = 0.002$~m hold
persistently (2\,\% of the initial deflection).
\item \textbf{Overshoot} $M_p$ --- maximum percentage excursion of $y$ beyond $y_\mathrm{ref}$
relative to the initial step amplitude.
\item \textbf{Steady-state error} $e_{ss}$ --- mean tracking error over the final 10\,\% of the
simulation horizon, measuring static accuracy.
\end{itemize}
 
\section{Simulations}
\label{sec:s}
\hl{Nie uzupelniac}


    
\setlength\tabcolsep{2pt}
\begin{table}
\caption{Example}
\label{table:qqq}
\centering
\begin{small}\scalebox{1}{
\begin{tabular}{lcccc}
\hline
\multicolumn{1}{c}{Algorithm\rule{0pt}{3mm}} & \multicolumn{1}{c}{$E^{\mathrm{pan}}_\mathrm{IAE}$} &\multicolumn{1}{c}{$E^{\mathrm{tilt}}_\mathrm{IAE}$} &\multicolumn{1}{c}{$E^{\mathrm{pan}}_\mathrm{MSE}$} &\multicolumn{1}{c}{$E^{\mathrm{tilt}}_\mathrm{MSE}$}\\
\hline
\rule{0pt}{3.5mm} PID &$3.58\times10^{-4} $ &$4.08\times10^{-4}$ &$4.18\times10^{-4}$ &$1.38\times10^{-3}$ \\
\rule{0pt}{3.5mm} DMC &$3.60\times10^{-4} $ &$ 4.08\times10^{-4}$ &$6.95\times10^{-7}$ &$9.22\times10^{-4}$ \\
\rule{0pt}{3.5mm} GPC &$1.24\times10^{-3} $ &$ 3.32\times10^{-4}$ &$6.56\times10^{-4}$ &$2.12\times10^{-2}$ \\
\hline
\end{tabular}}
\end{small}
\end{table}

    
	
	\begin{figure}[!htbp]
         \begin{center}
		\includegraphics[trim=0.0cm 0.1cm 0.0cm 0.0cm, clip=true, width=0.95\linewidth]{./Img/velma_b_v2.pdf}
		\end{center}
        \caption{Example}
        \label{fig:qqq}
    \end{figure}

    
	\section{Conclusions}
	\label{sec:Conclusions}
	\hl{Nie uzupelniac}
	
	\bibliographystyle{spmpsci}
	\bibliography{references}
	
	\section*{Statements and Declarations}
	
	The authors declare that no funds, grants, or other support were received during the preparation of this manuscript.
	
	The authors have no relevant financial or non-financial interests to disclose.
		
	The datasets generated during and/or analyzed during the current study are not publicly available due to complexity of the experiments, but are available from the corresponding author on reasonable request.
	
\end{document}




